% ============================================================================
% EXEMPLO FICTÍCIO — Relatório de Estágio como Prática Profissional
% Curso Técnico Integrado em Informática (PPC 2012) — IFRN Campus Ceará-Mirim
% ============================================================================
%
% Este arquivo é o irmão de "relatorio-metalinguistico.tex": mesmo modelo,
% mesma estrutura de seções, mas com um estágio fictício preenchido do
% início ao fim, para servir de referência visual de como o relatório
% pronto deve ficar.
%
% ATENÇÃO: dados, números, nomes de pessoas e de empresa deste exemplo são
% inventados para fins didáticos. Não citar nem reutilizar como se fossem
% reais.
%
% Compila com pdflatex (rode duas vezes por causa do sumário/referências).
% ============================================================================

\documentclass[article,12pt,a4paper]{abntex2}

\usepackage[T1]{fontenc}
\usepackage[utf8]{inputenc}
\usepackage[brazil]{babel}
\usepackage{times}
\usepackage{indentfirst}
\usepackage{graphicx}
\usepackage{microtype}

\renewcommand{\maketitlehookb}{}%% sem título em língua estrangeira (sem abstract)
\titulo{Relatório de Estágio Curricular no Setor de Suporte Técnico da Conecta Informática Ltda.}
\autor{Bruno Henrique Alves Farias \and Prof.\ Ricardo Nogueira de Souza\footnote{Professor orientador. Docente da área de Informática, IFRN Campus Ceará-Mirim.}}
\local{Ceará-Mirim -- RN}
\data{2026}
\instituicao{Instituto Federal de Educação, Ciência e Tecnologia do Rio Grande do Norte -- IFRN \\ Campus Ceará-Mirim}

\begin{document}
\maketitle

\begin{resumo}
O estágio curricular é um dos meios pelos quais o estudante do curso
Técnico Integrado em Informática aplica, em ambiente real de trabalho, os
conhecimentos desenvolvidos ao longo da formação técnica. Este relatório
tem como objetivo descrever e refletir sobre as atividades desenvolvidas
durante o estágio curricular supervisionado realizado no setor de suporte
técnico da empresa Conecta Informática Ltda., em Ceará-Mirim/RN, com carga
horária total de 300 horas, em jornada de 6 horas diárias. As atividades
desenvolvidas incluíram atendimento a chamados de suporte, instalação e
configuração de estações de trabalho, manutenção preventiva de
computadores e apoio à administração da rede local da empresa. Ao longo do
estágio, foram atendidos 87 chamados de suporte técnico e realizada a
configuração de 12 estações de trabalho novas. Conclui-se que o estágio
permitiu a aplicação prática de conteúdos estudados nas disciplinas de
Organização e Manutenção de Computadores e de Arquitetura de Redes de
Computadores, e contribuiu para o desenvolvimento de competências de
atendimento ao usuário não trabalhadas em sala de aula.

\vspace{\onelineskip}
\noindent
\textbf{Palavras-chave}: Estágio curricular. Suporte técnico. Manutenção de computadores. Redes de computadores.
\end{resumo}

\textual

\section{Introdução}

A busca por um estágio na área de suporte técnico surgiu do interesse em
aplicar, em ambiente real de trabalho, os conteúdos de manutenção de
computadores e redes estudados ao longo do curso Técnico Integrado em
Informática. A oportunidade de estágio na Conecta Informática Ltda.,
empresa de prestação de serviços de TI em Ceará-Mirim/RN, permitiu essa
aplicação prática, alinhada à área técnica do curso.

Este relatório tem como objetivo registrar e refletir sobre as atividades
desenvolvidas no setor de suporte técnico da empresa, ao longo do período
de estágio, em conformidade com o Plano de Estágio assinado pelo
estudante, pelo professor orientador e pelo supervisor técnico da
concedente.

O relatório está organizado da seguinte forma: a seção 2 apresenta os
dados formais do estágio; a seção 3 caracteriza a empresa concedente; a
seção 4 detalha as atividades desenvolvidas por tipo de tarefa; e a seção
5 traz as considerações finais.

\section{Estágio Curricular}

O estágio é regido pela Lei nº 11.788, de 25 de setembro de 2008 (Lei do
Estágio), e pela normativa interna do IFRN (Organização Didática, Arts.
307 a 319). Trata-se de estágio curricular não obrigatório, realizado a
partir do segundo ano do curso, conforme facultado pela Organização
Didática vigente.

O objetivo do estágio foi possibilitar ao estudante o exercício da prática
profissional, aliando teoria e prática, e facilitar seu ingresso no mundo
do trabalho, em conformidade com os objetivos previstos para o estágio
técnico na normativa do IFRN.

A carga horária total contratada foi de 300 horas, cumpridas em jornada de
6 horas diárias e 30 horas semanais, compatível com o horário das aulas do
estudante, conforme o Termo de Compromisso firmado entre o IFRN, a
concedente e o estudante. O estágio foi realizado no período de 3 de
março a 30 de junho de 2026.

\section{Caracterização da Concedente}

A Conecta Informática Ltda. é uma empresa de prestação de serviços de
tecnologia da informação, localizada no centro de Ceará-Mirim/RN, atuando
no atendimento a clientes pessoa física e pequenas empresas da região. O
fluxo de serviço da empresa consiste na abertura de chamados por telefone
ou aplicativo de mensagens, triagem pelo setor de suporte e atendimento
presencial ou remoto, conforme a complexidade do problema.

Os tipos de serviço prestados pela empresa incluem manutenção de
computadores e notebooks, configuração de redes domésticas e empresariais,
recuperação de dados e venda de equipamentos de informática.

O setor de suporte técnico, onde o estágio foi realizado, conta com 4
funcionários, incluindo o supervisor técnico responsável pela orientação
do estagiário, e dispõe de uma bancada de manutenção com ferramentas
próprias, um estoque de peças de reposição e uma rede local com 6
estações de trabalho para diagnóstico e testes.

\section{Atividades Desenvolvidas}

\subsection{Atendimento a chamados de suporte}

Foram atendidos, ao longo do estágio, 87 chamados de suporte técnico,
envolvendo problemas de lentidão, travamento, infecção por software
malicioso e falhas de conectividade de rede. Cada chamado era triado pelo
supervisor, e o estudante realizava o diagnóstico inicial e, sob
supervisão, a correção do problema. A atividade foi feita seguindo o
fluxo de atendimento da empresa, e permitiu desenvolver agilidade no
diagnóstico de problemas comuns de hardware e software, habilidade pouco
trabalhada nas aulas práticas do curso, que costumam envolver cenários já
previamente diagnosticados.

\subsection{Instalação e configuração de estações de trabalho}

Foram configuradas 12 estações de trabalho novas para clientes da empresa,
incluindo instalação de sistema operacional, drivers, pacote de
aplicativos de escritório e configuração de acesso à rede. Essa atividade
foi realizada seguindo um checklist padrão da empresa, e aplicou
diretamente os conteúdos de instalação e configuração de sistemas
estudados na disciplina de Organização e Manutenção de Computadores.

\subsection{Apoio à administração de rede}

O estagiário apoiou o supervisor técnico na configuração de três redes
locais em clientes empresariais, incluindo a configuração de roteadores,
criação de redes separadas para convidados e testes de velocidade de
conexão. A atividade foi feita sob supervisão direta, com o estagiário
executando as etapas técnicas sob orientação, e permitiu aplicar
conceitos de endereçamento IP e configuração de redes estudados na
disciplina de Arquitetura de Redes de Computadores.

\section{Considerações Finais}

O estágio atendeu às expectativas iniciais, e o tempo de 300 horas foi
suficiente para o estudante participar de um volume significativo de
atendimentos e desenvolver autonomia crescente ao longo do período. O
contato com os técnicos da empresa e com os clientes atendidos contribuiu
para o desenvolvimento de habilidades de comunicação e atendimento ao
público, complementares às habilidades técnicas.

Há uma correlação direta entre as atividades desenvolvidas no estágio e os
conteúdos das disciplinas de Organização e Manutenção de Computadores e de
Arquitetura de Redes de Computadores, o que confirma a unidade entre
teoria e prática exigida pela prática profissional. Como sugestão para
estagiários futuros na mesma área, recomenda-se buscar, desde o início,
anotar os chamados atendidos com detalhe, o que facilita bastante a
elaboração deste relatório ao final do estágio.

\postextual

\section*{Referências}

\newcommand{\referencia}[1]{\par\noindent\hangindent=1cm #1\par\medskip}

\referencia{BRASIL. \textit{Lei nº 11.788, de 25 de setembro de 2008}.
Dispõe sobre o estágio de estudantes. Brasília, DF: Diário Oficial da
União, 2008.}

\referencia{CONECTA INFORMÁTICA LTDA. \textit{Manual interno de
atendimento ao cliente}. Ceará-Mirim: Conecta Informática, 2024.}


\end{document}
